\subsection{Odabir teme i skupa podataka}\label{sec:podaci}
\textbf{Tema.} Predmet projekta je prepoznavanje prevarnih finansijskih transakcija na osnovu
karakteristika pojedinačnih transakcija. Zadatak se posmatra kao \textbf{nadzirana binarna
klasifikacija}, jer je za treniranje poznata ciljna klasa.

\textbf{Izvor.} Kao polazni skup koristi se \emph{Credit Card Transactions Fraud Detection Dataset}
sa platforme Kaggle \cite{kaggle_fraud}. U daljem radu analizira se stvarno korištena
verzija podataka. Ciljni atribut je \texttt{is\_fraud}: vrijednost 0 označava legitimnu,
a 1 prevarnu transakciju.

\begin{table}[htbp]\centering
\caption{Osnovne karakteristike korištenog skupa podataka}\label{tab:osnovni}
\begin{tabularx}{\linewidth}{@{}p{.44\linewidth}Y@{}}\toprule
Karakteristika & Podatak \\ \midrule
Naziv datoteke & \dopuni{naziv korištene datoteke} \\
Broj instanci & \dopuni{broj instanci} \\
Broj atributa & \dopuni{broj atributa} \\
Numerički atributi & \dopuni{broj i primjeri} \\
Kategorički atributi & \dopuni{broj i primjeri} \\
Ciljna varijabla & \texttt{is\_fraud} (0 -- legitimna; 1 -- prevarna) \\
Udio klase 1 & \dopuni{udio prevarnih transakcija u \%} \\
Način formiranja korištenog skupa & \dopuni{kratko opisati izbor redova i atributa} \\
\bottomrule\end{tabularx}
\end{table}

\textbf{Provjera uslova specifikacije.} Na korištenom skupu potvrditi najmanje 1000
instanci, najmanje 15 atributa, prisustvo numeričkih i kategoričkih atributa te
postojanje ciljne varijable. Navesti šta atributi predstavljaju i koje su eventualne
ograničenosti podataka. Ovdje ne treba prepisivati sve kolone originalnog skupa ako
nisu u stvarnoj verziji korištenoj za eksperimente.
\referencaKod{1.1 Odabir i učitavanje skupa podataka}
